\section{Rational curves on a contracted face}\label{sec:ruling}

For a genus-zero curve the scalar and balanced factors coincide. Write
$A_0(q)=\sum_{m\ge0}x_{0,m}q^m$. A family of resolved nodes shows how
a negative power of this factor can arise even though every cycle is
effective. We give the family first because it also computes the
contribution of a moving exceptional ruling.

Let $E$ be the nonsplit extension of $\cO$ by $\cO(-2)$ over $\PP^1$.
Then $E\simeq\cO(-1)^{\oplus2}$, since the boundary map from
$H^0(\cO)$ to $H^1(\cO(-2))$ is an isomorphism. Let
$u:\tot E\to\CC$ be the quotient coordinate. For a smooth complex
function $f$ on a closed disc, nonzero on the boundary, form the
fiber product $u=f(t)$ over its interior. A complex-horizontal
right inverse $h$ to $du$ gives this threefold an almost-complex
structure: in the splitting of a tangent vector as
$(a,k+h(df(a)))$ use multiplication by $i$ on $a$ and the vertical
complex structure on $k$. The projection to the disc is holomorphic
and each surface fiber retains its original complex structure.
For holomorphic $f$ this is the complex fiber product.

\begin{lemma}\label{lem:fold}
The series of compact cycles in this family, graded by their degree
over $\PP^1$, is $\Phi(A_0(q))^{\deg(f,0)}$ in the smooth algebra.
For holomorphic $f$ it is $A_0(q)^{\deg(f,0)}$ in $H$.
\end{lemma}
\begin{proof}
A compact irreducible curve projects constantly to the disc. The only
compact irreducible curve in $\tot E$ is its zero section, by negativity
after pullback to the normalization. Thus over a zero of $f$ there is
one rational curve, and there are no others. The full degree-$m$ cycle
locus is $\Sym^m(f^{-1}(0))$: the parametrization by sums is continuous
in support and current and is a continuous bijection from a compact
space to a Hausdorff space. The same construction over a compact path
without boundary zeros transports its unit. For separated zero sets
one can excise into disjoint targets, so their series multiply.

A positive simple zero can be deformed to the identity map, and there
the family is an open neighborhood of the zero section of $E$.
Its factor is $\Phi(A_0)$. At a negative regular zero the germ deforms
to complex conjugation by linearizing at the zero and moving within
the orientation-reversing linear isomorphisms. The local homotopies
avoid boundary zeros on a small disc. In the family
$f_s(x+iy)=x^2-s+iy$ there are no zeros for $s<0$ and two regular
zeros of opposite signs for $s>0$. The two factors therefore multiply
to one. Translation of any $f$ by a sufficiently small regular value
counts its simple zeros with Brouwer signs, proving the assertion.
For holomorphic $f$ the small translations are a holomorphic family
and the simple zeros are positive; the same computation is made
directly on complex cycles.
\end{proof}

For $f(t)=t^k-s$, continuation around $s=0$ permutes the $k$ supports
cyclically. The limiting single support has factor $A_0^k$. In bundle
coordinates its transition at $s=0$ is
$v'=z^2v+zt^k$, $t'=t$. Thus for distinct $k,l>r$ the $r$th
infinitesimal neighborhoods of the zero curves agree while their series
differ already in degree one. This is a dependence on the nonlinear
neighborhood that persists however large the prescribed jet order is.

\begin{proposition}\label{prop:ruled}
Let $p:X\to W$ be a birational contraction from a smooth projective
threefold with $K_X\simeq\cO_X$ to a projective variety. Suppose $p$ is an
isomorphism off a smooth divisor $D=\PP_B(V)$, and contracts exactly
its ruling fibers to the smooth connected curve $B\simeq p(D)$. If $b=g(B)$ and
$\ell$ is a fiber class, the complete universal series on the contracted
face is $A_0(Q^\ell)^{2b-2}$.
\end{proposition}
\begin{proof}
The pullback of an ample divisor on $W$ exposes the contracted face.
Every irreducible component there is a ruling fiber, and in class
$m\ell$ the cycle locus is $\Sym^m B$. All fibers have the same
integral class because $B$ is connected. An ample degree on $X$
distinguishes its multiples, even when $\ell$ is divisible in
$H_2(X;\ZZ)$. Deform $D$ to its normal
bundle in $X$ by blowing up $D\times0$ in $X\times\CC$ and deleting
the strict transform of the central fiber. The normal bundle is $K_D$.
A compact curve on the special fiber with contracted pushforward lies
over one point of $B$. On that ruling the normal bundle is $\cO(-2)$,
so the curve is the zero section. The relative fixed-label locus over
a closed parameter interval remains $\Sym^m B$ times that interval,
by the same compact continuous parametrization as in Lemma~\ref{lem:fold}.

Use the convention of lines in $V$. The relative Euler sequence gives
\[
 0\longrightarrow K_D\longrightarrow
 \cO_D(-1)\otimes\pi^*(V^\vee\otimes K_B)
 \xrightarrow{u}\pi^*K_B\longrightarrow0 ,
 \qquad \pi:D\to B.
\]
Replace the zero section on the right by a transverse smooth section
of $K_B$. Locally this is the family in Lemma~\ref{lem:fold}, since
on each ruling the Euler extension is nonsplit. A complex-linear
horizontal splitting gives the same fiberwise structures as in that
lemma and can be exchanged with a local product splitting by a homotopy.
Globally the projection to $B$ remains pseudoholomorphic. Each curve
of zero base degree therefore stays in one fiber; along the homotopy
of sections the relative symmetric powers of the compact zero loci
give compact support in every degree. The signed number of zeros of
the section is $2b-2$. Multiplying the local factors proves the assertion
after applying $\Phi$. Both complex series satisfy addition and
division in the integer label $m$, since every effective summand or
root of a contracted cycle is contracted. Lemma~\ref{lem:detect}
lifts the equality. For $b\ge1$ one can equally use a nonzero
holomorphic differential and count its zeros holomorphically.
The cotangent obstruction and its Euler number in the numerical
type III theory are due to Wilson \cite[\S4]{Wilson}.
\end{proof}

\section{Specialization across a conifold transition}\label{sec:transition}

Suppose a flat projective family over a smooth curve has smooth
algebraic total space of dimension four and central fiber with only
ordinary double points. Let $X_s$ be a smooth nearby fiber and let
$Y\to X_0$ be a projective small resolution. Assume the two smooth
threefolds have trivial canonical bundles and torsion-free second
integral homology. Fix the common-complement marking and a relatively
ample line bundle. Write $\pi:H_2(Y;\ZZ)\to H_2(X_s;\ZZ)$ for the
resulting quotient by the exceptional curve classes $[C_i]$.
Writing $U_X$ for the series in \eqref{eq:compact}, the universal
specialization identity of \cite[Theorem~3.2]{WueConifold} reads
\begin{equation}\label{eq:transition}
 U_{X_s}=\pi_*\left(U_Y\prod_i A_0(Q^{[C_i]})^{-1}\right).
\end{equation}
Here $\pi_*$ adds coefficients of classes with the same image.
Finiteness requires the exceptional normalization inside the parentheses.
In the blowup $\phi:\widetilde Y\to Y$ along the $C_i$, put
$j=\mathrm{PD}_{\widetilde Y}^{-1}\phi^*\mathrm{PD}_Y$. After normalization
all nonzero labels belong to
$F=\{\gamma\ne0:j\gamma\text{ is effective on }\widetilde Y\}$.
This set is closed under addition, avoids $\ker\pi$, and has finite
sublevel sets for the projected ample degree
\cite[Proposition~3.1]{WueConifold}. Indeed the ample degree of $j\gamma$
for a suitable polarization on the blowup is a fixed positive
multiple of the projected degree, since $j\gamma$ pairs to zero with
every exceptional divisor. The support assertion in \eqref{eq:transition}
in that generality uses the numerical reconstruction of
\cite[\S2]{WueConifold}. It assumes that the indicated connected unmarked
Gromov--Witten evaluations descend to the relevant relations in the
rational cycle algebra, including those arising from families of
complex targets. The analytic-family input to that descent is discussed
in \cite[Remark~1.8]{Pardon}. There is a direct argument using only
fixed-target counts when both endpoints also have
$H^1(\cO)=H^2(\cO)=0$.

The normalized series has integral balanced exponents
$\widetilde n_{\gamma,h}$, obtained by subtracting one at each
exceptional factor with its actual class. On the latter projective
domain let $G_h(q,u)$ denote the connected numerical series of the
elementary genus-$h$ curve. The signed-cluster comparison of
\cite[\S7.4]{WueProduct} identifies the connected invariants of each fixed
endpoint with the sum of $n_{\gamma,h}G_h(Q^\gamma,u)$. Write
$w=(2\sin(u/2))^2$; the coefficient of $q$ in $G_h$ is $w^{h-1}$.
The contracted face consists exactly of sums of the disjoint resolved
nodes, so subtracting their $G_0$ removes all connected contributions
on that face. For $\pi\gamma\ne0$ the ordinary degeneration of the
resolution gives its connected invariant as the empty-contact
relative invariant of $j\gamma$ on the blowup. The nonempty contact
strata have negative virtual dimension on the blowup side. In
particular every connected coefficient at $\gamma\notin F$ vanishes
after normalization \cite[\S2.1]{LLW}, \cite[\S3]{WueConifold}.

Induct on an ample degree on the resolution. If $\gamma\notin F$,
every proper root of $\gamma$ is outside $F$, since $F$ is closed
under multiplication by positive integers. After the lower-degree
terms have vanished, the connected coefficient at $\gamma$ is
$\sum_h\widetilde n_{\gamma,h}w^{h-1}=0$. This is a finite Laurent
polynomial and its least power has a distinct leading order in $u$.
All its coefficients therefore vanish. Thus every remaining exponent
is supported in $F$, and its product may be pushed forward in finite
sums. The connected smoothing formula in the total-space class adds
the fixed-target resolution counts with that image, also for classes
ineffective on the smoothing \cite[\S3, proof of Theorem~3.2]{WueConifold},
\cite[Theorem~2.16 and \S3]{LW}. Comparing with the pushed-forward
product, the same induction now in the smoothing degree identifies
its exponents. This proves \eqref{eq:transition} on the stated
projective domain without a theory of evaluations on general families.
With the numerical reconstruction assumed in the broader case, the
support theorem in \cite[\S3]{WueConifold} and successive division by
factors in $F$ give the same exponent formula. In either case,
for every nonzero effective $\beta$ on the smoothing,
\[
 n_{\beta,h}(X_s)=\sum_{\substack{\gamma\in F\\\pi\gamma=\beta}}
                       \widetilde n_{\gamma,h}.
\]
The proof applies in the monoid generated by the effective smoothing
classes and $\pi(F)$, extending the series by zero. This also describes
the cancellation when an image class is ineffective. Two small
resolutions have the same blowup and their normalized series agree
under its common zero-contact identification \cite[\S3]{WueConifold}.
The same summation for scalar product exponents preserves their
integrality whenever the source exponents are integral
\cite[Corollary~3.3]{WueConifold}. The scalar exponent in
\S\ref{sec:intro} explains the need for this hypothesis.

One can bound the sum explicitly on the classical transition of a
quintic containing a plane \cite{Rossi}.
Let $f,g$ be general quartics on $\PP^4$. The central quintic
$x_3f+x_4g=0$ has sixteen nodes on the plane $x_3=x_4=0$.
Indeed the two restricted quartics meet transversely, and at each
intersection $x_3,x_4,f,g$ are local coordinates. Away from the
plane smoothness follows from Bertini. In $\PP^1_{[s:t]}\times\PP^4$
the equations $s x_3+t x_4=0$ and $-s g+t f=0$ give its projective
small resolution $Y$. The projection selects the kernel line of the
two-by-two coefficient matrix, with an entire ruling over each node.
Adjunction gives $K_Y=\cO_Y$. Adding a general multiple of a quintic
nonvanishing at all sixteen nodes gives a smoothing with smooth total
space near the central fiber. Write $B,H$ for the hyperplane classes
from the two factors. Integral Lefschetz for the two ample divisors gives
coordinates $(a,d)=(B\cdot\gamma,H\cdot\gamma)$ on $H_2(Y;\ZZ)$,
and the smoothing map is $(a,d)\mapsto d$. Each exceptional curve
has class $(1,0)$.

In this example $F$ satisfies $0\le a\le 4d$, $d\ge1$. To see the
upper bound, write $u=f/s=g/t$ and $v=x_4/s=-x_3/t$. These sections
have classes $4H-B$ and $H-B$. The sections $u,vh$, with $h$ ranging
over cubics, generate the ideal of the exceptional curves after twisting
by $4H-B$. Indeed their common ideal is locally $(u,v)$, the two normal
coordinates in the small resolution of the node. Thus on the blowup
$\phi^*(4H-B)-E$ is generated by global sections, where $E$ is the
sum of the exceptional divisors. Pairing this nef divisor and $\phi^*B$
with $j\gamma$ proves both inequalities. A line in the plane missing
the nodes lifts to the graph of the degree-four map $[f:g]$ on the
line. Its proper transform has class $j(4,1)$, so the upper bound is
attained in $F$. After division by $A_0(q)^{16}$ the coefficient in
smoothing degree $d>0$ is therefore a polynomial in the exceptional
variable $q$ of degree at most $4d$. Membership in this support bound
alone does not establish that a coefficient is nonzero.

In the four-node toric transition let $\beta=(1,1,0)$ be the
primitive line class of the smoothing plane. In the resolution basis
of \cite[\S4]{WueConifold} the marking takes
$(n_1,\ldots,n_6)$ to $(n_2,n_3,n_4)$. The following individual source and smoothing coefficients use
the stable-pair evaluation of geometric local classes assumed in
\cite[\S\S5--6]{WueConifold}; see also \cite[Remark~1.8]{Pardon}.
The conifold specialization proved above uses the fixed-target
Gromov--Witten comparison. Of the 37 candidate lifts of
$\beta$, eight contribute. Their complete coefficients are given in
\cite[(21)--(24)]{WueConifold}. The reconstruction also uses a bound on the Euler characteristic of a pure curve. The quadratic terms cancel upon specialization, so
the eight polynomials sum to $3x_{0,1}$. On the smoothing the minimal
pair spaces for lines and conics are the dual plane and $\PP^5$.
Their signed Euler characteristics $3$ and $-6$ give
\[
 U_{X_s,\beta}=3x_{0,1},\qquad
 U_{X_s,2\beta}=3x_{0,2}+3x_{0,1}^2-6x_{0,1}
 \qquad\text{\cite[Proposition~6.1]{WueConifold}}.
\]
The line spans an extremal ray and the only effective splitting of its
double is $\beta+\beta$. Thus the first factor on the ray is
$A_0(Q^\beta)^3$. At $2\beta$ its contribution is
$3x_{0,2}+3x_{0,1}^2$. The remaining $-6x_{0,1}$ is the leading
coefficient of $A_0(Q^{2\beta})^{-6}$.
